\documentclass{beamer}

%lembre-se de configurar o pdflatex com:
%pdflatex -shell-escape -synctex=1 -interaction=nonstopmode  -output-directory=build  %.tex

\usepackage{aula}

\title{Aula modelo}
\date{\today}

\begin{document}
    
\input{capa}

\sumario

\section{SQLAlchemy}

% =====================================================================
% SLIDE 1
% =====================================================================

% =====================================================================
% SLIDE 1
% =====================================================================
\begin{frame}{ORM: Object Relational Mapping}

\textbf{ORM} (Object Relational Mapping) é uma técnica para mapear tabelas do banco de dados para classes Python.

\vspace{0.4cm}

\textbf{Sem ORM}

\begin{itemize}
    \item Escrever SQL manualmente
    \item Converter resultados para objetos
\end{itemize}

\vspace{0.3cm}

\textbf{Com ORM}

\begin{itemize}
    \item Classes representam tabelas
    \item Objetos representam registros
    \item Métodos substituem comandos SQL simples
\end{itemize}

\vspace{0.4cm}

\textbf{Benefícios}

\begin{itemize}
    \item Menos SQL repetitivo
    \item Código mais organizado
    \item Compatibilidade entre bancos
\end{itemize}

\end{frame}

% =====================================================================
% SLIDE 2
% =====================================================================
\begin{frame}[fragile]{SQLAlchemy: Conectando ao Banco}

\textbf{GeoPackage (SQLite)}

\begin{minted}{python}
from sqlalchemy import create_engine

engine = create_engine(
    "sqlite:///dados.gpkg"
)
\end{minted}

\vspace{0.3cm}

\textbf{PostGIS (PostgreSQL)}

\begin{minted}{python}
from sqlalchemy import create_engine

engine = create_engine(
    "postgresql+psycopg://"
    "postgres:postgres"
    "@localhost/geodb"
)
\end{minted}

\vspace{0.3cm}

O restante do código ORM é praticamente o mesmo para ambos.

\end{frame}

% =====================================================================
% SLIDE 3
% =====================================================================
\begin{frame}[fragile]{Mapeando uma Tabela para uma Classe}

\begin{minted}{python}
from sqlalchemy.orm import (
    DeclarativeBase,
    Mapped,
    mapped_column
)

class Base(DeclarativeBase):
    pass

class Cidade(Base):

    __tablename__ = "cidades"

    id: Mapped[int] = mapped_column(
        primary_key=True
    )

    nome: Mapped[str]

    populacao: Mapped[int]
\end{minted}

\vspace{0.2cm}

\textbf{Correspondência}

\begin{itemize}
    \item Classe → tabela
    \item Atributo → coluna
    \item Objeto → registro
\end{itemize}

\end{frame}

% =====================================================================
% SLIDE 4
% =====================================================================
\begin{frame}[fragile]{Criando Tabelas}

\begin{minted}{python}
Base.metadata.create_all(
    engine
)
\end{minted}

\vspace{0.3cm}

SQL equivalente:

\begin{minted}{sql}
CREATE TABLE cidades (
    id INTEGER PRIMARY KEY,
    nome TEXT,
    populacao INTEGER
);
\end{minted}

\vspace{0.3cm}

O ORM gera a estrutura automaticamente.

\end{frame}

% =====================================================================
% SLIDE 5
% =====================================================================
\begin{frame}[fragile]{Inserindo Registros}

\begin{minted}{python}
from sqlalchemy.orm import Session

with Session(engine) as session:

    cidade = Cidade(
        nome="Curitiba",
        populacao=1774000
    )

    session.add(cidade)

    session.commit()
\end{minted}

\vspace{0.3cm}

Equivalente SQL:

\begin{minted}{sql}
INSERT INTO cidades
(nome, populacao)
VALUES
('Curitiba', 1774000);
\end{minted}

\end{frame}

% =====================================================================
% SLIDE 6
% =====================================================================
\begin{frame}[fragile]{Consultando Registros}

\begin{minted}{python}
from sqlalchemy import select

with Session(engine) as session:

    stmt = select(Cidade)

    resultado = session.scalars(
        stmt
    )

    for cidade in resultado:
        print(cidade.nome)
\end{minted}

\vspace{0.3cm}

Equivalente SQL:

\begin{minted}{sql}
SELECT *
FROM cidades;
\end{minted}

\end{frame}

% =====================================================================
% SLIDE 7
% =====================================================================
\begin{frame}[fragile]{Filtros e Ordenação}

\begin{minted}{python}
from sqlalchemy import select

stmt = (
    select(Cidade)
    .where(
        Cidade.populacao > 1000000
    )
    .order_by(
        Cidade.nome
    )
)
\end{minted}

\vspace{0.3cm}

Equivalente SQL:

\begin{minted}{sql}
SELECT *
FROM cidades
WHERE populacao > 1000000
ORDER BY nome;
\end{minted}

\end{frame}

% =====================================================================
% SLIDE 8
% =====================================================================
\begin{frame}[fragile]{Geometrias com PostGIS e GeoPackage}

\textbf{GeoAlchemy2}

\begin{minted}{python}
from geoalchemy2 import Geometry

class Cidade(Base):

    __tablename__ = "cidades"

    id = mapped_column(
        primary_key=True
    )

    nome = mapped_column()

    geom = mapped_column(
        Geometry(
            "POINT",
            srid=4326
        )
    )
\end{minted}

\vspace{0.3cm}

A coluna geom passa a ser tratada como uma geometria espacial pelo ORM.

\end{frame}

% =====================================================================
% SLIDE 9
% =====================================================================
\begin{frame}{ORM em Aplicações GIS}

\textbf{GeoPackage}

\begin{itemize}
    \item Aplicações desktop
    \item Dados locais
    \item Arquivo único
\end{itemize}

\vspace{0.3cm}

\textbf{PostGIS}

\begin{itemize}
    \item Múltiplos usuários
    \item Grandes volumes de dados
    \item Serviços web e APIs
\end{itemize}

\vspace{0.4cm}

\textbf{SQLAlchemy + GeoAlchemy2}

\begin{itemize}
    \item Abstração do banco
    \item Código reutilizável
    \item Integração com FastAPI, Flask e GeoPandas
\end{itemize}

\end{frame}

\section{SQLAlchemy com FastAPI}
\begin{frame}{ORM + FastAPI}

\textbf{FastAPI} é um framework moderno para construção de APIs REST em Python.

\vspace{0.4cm}

\textbf{Integração típica}

\begin{center}
Cliente
$\rightarrow$
FastAPI
$\rightarrow$
SQLAlchemy
$\rightarrow$
PostgreSQL/PostGIS
\end{center}

\vspace{0.4cm}

\textbf{Vantagens}

\begin{itemize}
    \item APIs REST rapidamente
    \item Integração nativa com JSON
    \item Documentação automática (Swagger)
    \item Suporte a PostgreSQL e PostGIS
\end{itemize}

\end{frame}

% =====================================================================
% SLIDE 2
% =====================================================================
\begin{frame}[fragile]{Configuração da Conexão}

\begin{minted}{python}
from sqlalchemy import create_engine

DATABASE_URL = (
    "postgresql+psycopg://"
    "postgres:postgres"
    "@localhost/geodb"
)

engine = create_engine(
    DATABASE_URL
)
\end{minted}

\vspace{0.4cm}

\textbf{Sessões do banco}

\begin{minted}{python}
from sqlalchemy.orm import Session

with Session(engine) as session:
    ...
\end{minted}

\end{frame}

% =====================================================================
% SLIDE 3
% =====================================================================
\begin{frame}[fragile]{Modelo ORM}

\begin{minted}{python}
from sqlalchemy.orm import (
    DeclarativeBase,
    mapped_column
)

class Base(DeclarativeBase):
    pass

class Cidade(Base):

    __tablename__ = "cidades"

    id = mapped_column(
        primary_key=True
    )

    nome = mapped_column()

    populacao = mapped_column()
\end{minted}

\vspace{0.2cm}

A classe Cidade representa a tabela do banco.

\end{frame}

% =====================================================================
% SLIDE 4
% =====================================================================
\begin{frame}[fragile]{Primeira Rota REST}

\begin{minted}{python}
from fastapi import FastAPI

app = FastAPI()

@app.get("/")
def raiz():

    return {
        "mensagem":
        "API funcionando"
    }
\end{minted}

\vspace{0.3cm}

Executar:

\begin{minted}{bash}
uvicorn app:app --reload
\end{minted}

\end{frame}

% =====================================================================
% SLIDE 5
% =====================================================================
\begin{frame}[fragile]{Consultando Dados}

\begin{minted}{python}
from sqlalchemy import select
from sqlalchemy.orm import Session

@app.get("/cidades")

def listar_cidades():

    with Session(engine) as session:

        stmt = select(Cidade)

        cidades = session.scalars(
            stmt
        ).all()

        return cidades
\end{minted}

\vspace{0.2cm}

Endpoint:

\begin{minted}{text}
/cidades
\end{minted}

\end{frame}

% =====================================================================
% SLIDE 6
% =====================================================================
\begin{frame}[fragile]{Inserindo Registros}

\begin{minted}{python}
@app.post("/cidades")

def criar_cidade():

    with Session(engine) as session:

        cidade = Cidade(
            nome="Campinas",
            populacao=1214000
        )

        session.add(cidade)

        session.commit()

        return {"status":"ok"}
\end{minted}

\end{frame}

% =====================================================================
% SLIDE 7
% =====================================================================
\begin{frame}[fragile]{Recebendo JSON}

\begin{minted}{python}
from pydantic import BaseModel

class CidadeDTO(BaseModel):

    nome: str
    populacao: int


@app.post("/cidades")

def criar(cidade: CidadeDTO):

    print(cidade.nome)

    return {"status":"ok"}
\end{minted}

\vspace{0.3cm}

Exemplo JSON:

\begin{minted}{json}
{
  "nome":"Curitiba",
  "populacao":1774000
}
\end{minted}

\end{frame}

% =====================================================================
% SLIDE 8
% =====================================================================
\begin{frame}[fragile]{Consultando por ID}

\begin{minted}{python}
@app.get("/cidades/{id}")

def buscar(id: int):

    with Session(engine) as session:

        cidade = session.get(
            Cidade,
            id
        )

        return cidade
\end{minted}

\vspace{0.3cm}

Exemplos:

\begin{minted}{text}
/cidades/1
/cidades/2
\end{minted}

\end{frame}

% =====================================================================
% SLIDE 9
% =====================================================================
\begin{frame}[fragile]{Swagger Automático}

\begin{itemize}
    \item Documentação gerada automaticamente
    \item Testes de API diretamente no navegador
    \item Visualização dos modelos Pydantic
\end{itemize}

\vspace{0.3cm}

Após iniciar a aplicação:

\begin{minted}{text}
http://localhost:8000/docs
\end{minted}

\vspace{0.3cm}

Permite executar:

\begin{itemize}
    \item GET
    \item POST
    \item PUT
    \item DELETE
\end{itemize}

Sem escrever código cliente.

\end{frame}

% =====================================================================
% SLIDE 10
% =====================================================================
\begin{frame}{Arquitetura Completa}

\begin{center}

Cliente (Browser, Python, QGIS)

\vspace{0.2cm}

$\downarrow$

\vspace{0.2cm}

FastAPI

\vspace{0.2cm}

$\downarrow$

\vspace{0.2cm}

SQLAlchemy ORM

\vspace{0.2cm}

$\downarrow$

\vspace{0.2cm}

PostgreSQL / PostGIS

\end{center}

\vspace{0.5cm}

\textbf{Benefícios}

\begin{itemize}
    \item Código organizado
    \item Menos SQL manual
    \item APIs REST modernas
    \item Integração com sistemas GIS
\end{itemize}

\end{frame}

\end{document}
